\documentclass[12pt, a4paper]{article}

% ── Packages ─────────────────────────────────────────────────────────────────
\usepackage[a4paper, left=2.5cm, right=2.5cm, top=3cm, bottom=2.5cm, headheight=15pt]{geometry}
\usepackage{fontenc}
\usepackage{inputenc}
\usepackage{microtype}
\usepackage{setspace}
\usepackage{parskip}
\usepackage{titlesec}
\usepackage{titletoc}
\usepackage{fancyhdr}
\usepackage{booktabs}
\usepackage{array}
\usepackage{longtable}
\usepackage{tabularx}
\usepackage{amsmath}
\usepackage{amssymb}
\usepackage{hyperref}
\usepackage{enumitem}
\usepackage{mdframed}
\usepackage{caption}
\usepackage{float}
\usepackage{ragged2e}
\usepackage{listings}
\usepackage{xcolor}

% ── Hyperref setup ───────────────────────────────────────────────────────────
\hypersetup{
    colorlinks=false,
    hidelinks,
    pdftitle={Qwen3-8B Fine-Tuning for Structured Tool Use},
    pdfauthor={Dhruva N},
    pdfsubject={Project Report},
}

% ── Code listing style ───────────────────────────────────────────────────────
\lstset{
    basicstyle=\ttfamily\small,
    backgroundcolor=\color{gray!8},
    frame=single,
    framerule=0.4pt,
    breaklines=true,
    columns=fullflexible,
    keepspaces=true,
}

% ── Section title formatting ─────────────────────────────────────────────────
\titleformat{\section}
  {\Large\bfseries}
  {\thesection.}{0.6em}{}
  [\vspace{-4pt}\rule{\linewidth}{0.5pt}\vspace{2pt}]

\titleformat{\subsection}
  {\normalsize\bfseries}
  {\thesubsection}{0.5em}{}

\titleformat{\subsubsection}
  {\normalsize\bfseries\itshape}
  {\thesubsubsection}{0.5em}{}

\titlespacing{\section}{0pt}{22pt}{8pt}
\titlespacing{\subsection}{0pt}{14pt}{5pt}
\titlespacing{\subsubsection}{0pt}{10pt}{4pt}

% ── Header / Footer ──────────────────────────────────────────────────────────
\pagestyle{fancy}
\fancyhf{}
\fancyhead[L]{}
\fancyhead[R]{}
\fancyfoot[L]{}
\fancyfoot[R]{\small\thepage}
\renewcommand{\headrulewidth}{0pt}
\renewcommand{\footrulewidth}{0pt}

% ── Abstract box ─────────────────────────────────────────────────────────────
\newmdenv[
  linewidth=0pt,
  innerleftmargin=0pt,
  innerrightmargin=0pt,
  innertopmargin=4pt,
  innerbottommargin=4pt,
  skipabove=8pt,
  skipbelow=8pt
]{abstractbox}

\newmdenv[
  linewidth=0.6pt,
  innerleftmargin=12pt,
  innerrightmargin=12pt,
  innertopmargin=8pt,
  innerbottommargin=8pt,
  skipabove=6pt,
  skipbelow=6pt
]{equationbox}

% ── Table helpers ─────────────────────────────────────────────────────────────
\newcommand{\thead}[1]{\textbf{#1}}
\newcolumntype{L}[1]{>{\RaggedRight\arraybackslash}p{#1}}
\newcolumntype{C}[1]{>{\Centering\arraybackslash}p{#1}}

% ── Misc ─────────────────────────────────────────────────────────────────────
\setlength{\parindent}{0pt}
\setlength{\parskip}{6pt}
\onehalfspacing
\captionsetup{font=small, labelfont=bf, textfont=it, skip=4pt}
\renewcommand{\arraystretch}{1.3}

% ─────────────────────────────────────────────────────────────────────────────
\begin{document}

% ── COVER PAGE ───────────────────────────────────────────────────────────────
\thispagestyle{empty}
\singlespacing
\vspace*{0.6cm}

\begin{center}
{\small\textbf{\MakeUppercase{Project Report}}}\\[6pt]
{\rule{\linewidth}{1pt}}\\[8pt]

{\LARGE\bfseries
Fine-Tuning Qwen3-8B for\\[4pt]
Structured Tool Use
}\\[8pt]
{\rule{\linewidth}{1pt}}\\[10pt]

{\large\textit{SFT and GRPO Pipeline with Tinker Remote Infrastructure}}\\[18pt]

{\large\textbf{Dhruva N}}\\[4pt]
{Department of Data Science and Artificial Intelligence}\\[4pt]
{April 2026}
\end{center}

\vspace{0.7cm}

\begin{abstractbox}
\textbf{Abstract}\\[6pt]
This report documents the design, implementation, and results of a fine-tuning pipeline
for \texttt{Qwen/Qwen3-8B}, targeting reliable multi-step structured tool calling. The
pipeline proceeds in three stages: zero-shot baseline evaluation, Supervised Fine-Tuning
(SFT) with LoRA adapters, and Group Relative Policy Optimization (GRPO) applied on top
of the SFT checkpoint. Training is performed on Tinker remote GPUs; local dry-run
validation uses a lightweight surrogate model (\texttt{sshleifer/tiny-gpt2}) to verify
pipeline wiring without consuming remote compute. The training dataset is entirely synthetic,
generated from a library of 18 tool templates covering single-step and multi-step tool
chains. Five programmatic reward signals drive GRPO without requiring any human-annotated
preference pairs. Evaluation compares all three stages across tool selection accuracy,
argument accuracy, schema compliance, and multi-step chain success, with latency measurements
included. GRPO produces the best argument accuracy (0.321) and the highest multi-step success
rate (0.413), while also reducing the fraction of zero-call failures in multi-step scenarios.
The end-to-end pipeline is fully reproducible via a single shell entrypoint.
\end{abstractbox}

\vspace{0.5cm}

\begin{center}
\small
\textbf{Keywords:} Qwen3-8B $\cdot$ LLM Fine-Tuning $\cdot$ SFT $\cdot$ GRPO $\cdot$ LoRA $\cdot$
Tool Calling $\cdot$ Function Calling $\cdot$ Synthetic Data $\cdot$ Tinker $\cdot$
Multi-Step Reasoning $\cdot$ Reward Functions $\cdot$ W\&B
\end{center}

\newpage
\onehalfspacing

% ── TABLE OF CONTENTS ────────────────────────────────────────────────────────
\tableofcontents
\newpage

% =============================================================================
\section{Project Overview}
% =============================================================================

This repository implements a four-stage pipeline for fine-tuning
\texttt{Qwen/Qwen3-8B} to reliably produce structured tool calls in response
to natural language queries. The central research question is whether a model
that has not been specifically instruction-tuned for a tool-calling schema can
be made to emit well-formed, executable tool calls through SFT alone, or
whether a reinforcement learning stage (GRPO) is needed to enforce compliance
with the output schema and improve multi-step chain quality.

The source repository for this report is:
\href{https://github.com/dhruvanmurthy/Qwen3-8B-FineTuning}{GitHub repository: Qwen3-8B-FineTuning}.

\subsection{Motivation}

Large language models are increasingly used as the reasoning backbone of
agentic systems. In such architectures the model must not merely answer
questions in natural language --- it must decide which external tool (API,
calculator, calendar service, etc.) to invoke, supply the correct arguments,
and, in multi-step scenarios, chain several tool calls in the right order
before synthesising a final response. This is demanding because:

\begin{itemize}[leftmargin=1.5em, itemsep=4pt]
  \item The model must respect a strict output schema even when its pre-training
    distribution favours natural language responses.
  \item In multi-step tasks the model must keep track of which tools have
    already been called and what their results were.
  \item Argument values must be grounded in the user query, not hallucinated.
\end{itemize}

Supervised fine-tuning is a natural first step, but imitation learning on
demonstration data does not provide a direct gradient signal towards schema
compliance. GRPO, which assigns relative rewards across a group of sampled
completions, supplies exactly that signal --- and does so without a separately
trained reward model or human preference labels.

\subsection{Pipeline Stages}

The pipeline is organised into four stages executed in sequence:

\begin{enumerate}[leftmargin=1.5em, itemsep=5pt]
  \item \textbf{Stage 0 -- Baseline Evaluation.}
    Evaluate zero-shot performance of the base \texttt{Qwen/Qwen3-8B} model
    with no fine-tuning. Establishes the floor for all comparisons.
  \item \textbf{Stage 1 -- SFT.}
    Fine-tune LoRA adapters on synthetic tool-use conversations using
    cross-entropy (next-token prediction) loss. Uses Tinker remote GPUs.
  \item \textbf{Stage 2 -- GRPO.}
    Resume training from the SFT checkpoint using Group Relative Policy
    Optimization with five programmatic reward signals. Same LoRA rank as SFT.
  \item \textbf{Stage 3 -- Comparison.}
    Run the full evaluation suite against all three checkpoints and produce
    a side-by-side metrics table.
\end{enumerate}

\subsection{Repository Layout}

\begin{table}[H]
\centering
\caption{Key files and directories in the repository.}
\begin{tabularx}{\linewidth}{L{6.2cm} L{8.8cm}}
\toprule
\thead{Path} & \thead{Purpose} \\
\midrule
\texttt{src/train.py} & SFT training script (Tinker-based) \\
\texttt{src/train\_grpo.py} & GRPO training script (Tinker-based) \\
\texttt{src/evaluate.py} & Three-stage evaluation harness \\
\texttt{src/rewards.py} & Programmatic reward functions \\
\texttt{src/data\_loader.py} & Dataset loading and preprocessing \\
\texttt{src/constants.py} & Shared constants including the system prompt \\
\texttt{scripts/run\_pipeline.sh} & Canonical end-to-end entrypoint \\
\texttt{scripts/generate\_synthetic.py} & Synthetic dataset generator \\
\texttt{scripts/prepare\_datasets.sh} & Data generation and tokenization \\
\texttt{configs/dataset\_config.yaml} & Dataset source and split configuration \\
\texttt{data/raw/synthetic/} & Generated JSONL training examples \\
\texttt{data/processed/} & Tokenized Arrow splits and test raw JSONL \\
\texttt{outputs/} & Evaluation JSON files and training artifacts \\
\texttt{wandb/} & W\&B run logs from dry-run validation sessions \\
\bottomrule
\end{tabularx}
\end{table}

\newpage

% =============================================================================
\section{Base Model and Infrastructure}
% =============================================================================

\subsection{Base Model: Qwen3-8B}

The base model for all experiments is \texttt{Qwen/Qwen3-8B}, the 8-billion
parameter instruct variant from Alibaba Cloud's Qwen3 series. Key properties
relevant to this project:

\begin{itemize}[leftmargin=1.5em, itemsep=4pt]
  \item Trained on a large multilingual corpus with instruction tuning and
    RLHF alignment.
  \item Supports extended context lengths up to 32\,768 tokens (pipeline uses
    at most 2\,048 tokens).
  \item Includes a built-in thinking mode; the reward extraction code in
    \texttt{src/rewards.py} strips \texttt{<think>} blocks before parsing tool
    calls to avoid false negatives.
  \item Achieves strong zero-shot tool-selection performance
    (\textasciitilde78\% on the synthetic benchmark), providing a meaningful
    baseline above which both fine-tuning stages are measured.
\end{itemize}

\subsection{Training Infrastructure: Tinker}

All SFT and GRPO training passes run on \textbf{Tinker} remote GPUs, a managed
machine learning infrastructure service. The local machine is responsible only
for data preparation and reward grading; forward/backward passes and parameter
updates are executed remotely. This architecture has two practical
consequences:

\begin{itemize}[leftmargin=1.5em, itemsep=4pt]
  \item \textbf{No local GPU required for full training.} The pipeline requires
    only \texttt{TINKER\_API\_KEY} in the environment.
  \item \textbf{Local dry-run mode.} Setting \texttt{LOCAL\_VALIDATE=true}
    swaps the base model to \texttt{sshleifer/tiny-gpt2} and enables
    \texttt{--dry-run} on all training scripts, allowing complete end-to-end
    pipeline wiring to be tested without Tinker API access or remote compute.
\end{itemize}

\subsection{Local Validation Environment}

Local dry-run validation was performed using \texttt{sshleifer/tiny-gpt2} as a
surrogate model to verify end-to-end pipeline wiring without consuming remote
compute. Key environment details:

\begin{table}[H]
\centering
\caption{Local validation environment from W\&B metadata.}
\begin{tabularx}{\linewidth}{L{3.5cm} L{10.7cm}}
\toprule
\thead{Property} & \thead{Value} \\
\midrule
OS & Linux (WSL2, kernel 6.6.87.2-microsoft-standard) \\
Python & CPython 3.11.15 \\
GPU & NVIDIA GeForce RTX 3060 Laptop GPU (6\,GB VRAM) \\
CUDA & 13.0 \\
W\&B CLI & 0.25.1 \\
Host & DHRUVA-LAPTOP \\
Conda env & rl-lab \\
\bottomrule
\end{tabularx}
\end{table}

Local dry-run executions confirmed that the SFT and GRPO scripts reach
checkpointing and reward-logging code paths correctly, and that the pipeline
completes without errors under \texttt{LOCAL\_VALIDATE=true}.

\subsection{Experiment Tracking: Weights \& Biases}

W\&B project \texttt{qwen3-8b-tool-use} is used for experiment tracking.
The SFT script logs \texttt{train/epoch}, \texttt{train/loss}, and
\texttt{train/samples\_seen} at each step. The GRPO script logs
\texttt{train/reward\_mean}, \texttt{train/frac\_degenerate},
\texttt{train/n\_datums}, and \texttt{train/n\_degenerate\_groups}.
W\&B is initialised in \textit{disabled} mode when \texttt{WANDB\_API\_KEY}
is absent, so all pipeline runs succeed silently without it.

\newpage

% =============================================================================
\section{Dataset Design and Preparation}
% =============================================================================

\subsection{Synthetic Data Strategy}

The dataset used for training and evaluation is \textbf{entirely synthetic},
generated by \texttt{scripts/generate\_synthetic.py}. This choice was made
deliberately:

\begin{itemize}[leftmargin=1.5em, itemsep=4pt]
  \item No licensing constraints from external corpora.
  \item Full control over tool diversity, argument type coverage, and
    multi-step chain structure.
  \item Deterministic reproduction with a fixed seed (\texttt{seed=42}).
\end{itemize}

The generator constructs two JSONL files:

\begin{itemize}[leftmargin=1.5em, itemsep=3pt]
  \item \texttt{data/raw/synthetic/synthetic\_single.jsonl} --- single-tool-call examples.
  \item \texttt{data/raw/synthetic/synthetic\_multistep.jsonl} --- multi-tool-chain examples.
\end{itemize}

\subsection{Tool Catalogue}

The generator defines a catalogue of 18 tools. Table~\ref{tab:tools} lists
the tools used for single-step generation. Multi-step chains compose subsets
of these tools.

\begin{table}[H]
\centering
\caption{Synthetic tool catalogue (single-step generators).}
\label{tab:tools}
\begin{tabularx}{\linewidth}{L{4.2cm} L{9.8cm}}
\toprule
\thead{Tool Name} & \thead{Description / Parameters} \\
\midrule
\texttt{get\_weather} & Current weather for a city; \texttt{city} (str, required), \texttt{units} (C/F, optional) \\
\texttt{get\_stock\_price} & Stock price for a ticker; \texttt{symbol} (str), \texttt{currency} (optional) \\
\texttt{search\_web} & Web search; \texttt{query} (str), \texttt{num\_results} (int, optional) \\
\texttt{send\_email} & Send email; \texttt{to}, \texttt{subject}, \texttt{body} (str, required), \texttt{cc} (optional) \\
\texttt{create\_calendar\_event} & Calendar entry; \texttt{title}, \texttt{date}, \texttt{time}, \texttt{duration\_minutes}, \texttt{attendees} \\
\texttt{translate\_text} & Translation; \texttt{text}, \texttt{target\_lang} (required), \texttt{source\_lang} (optional) \\
\texttt{get\_news} & News articles; \texttt{topic}, \texttt{num\_articles}, \texttt{language} \\
\texttt{calculate} & Math expression; \texttt{expression} (str, required) \\
\texttt{get\_restaurant\_info} & Restaurant lookup; \texttt{name} (str, required) \\
\texttt{convert\_currency} & Currency conversion; \texttt{amount}, \texttt{from\_currency}, \texttt{to\_currency} \\
\texttt{convert\_units} & Unit conversion; \texttt{value}, \texttt{from\_unit}, \texttt{to\_unit} \\
\texttt{set\_reminder} & Reminder; \texttt{message}, \texttt{datetime} \\
\bottomrule
\end{tabularx}
\end{table}

Multi-step categories include: translate+weather, stock+convert, weather+remind,
search+email, and other combinations, each exercising 2--3 tool calls in a
single training example.

\subsection{Data Format}

Every generated example carries the following structured fields:

\begin{itemize}[leftmargin=1.5em, itemsep=3pt]
  \item \textbf{\texttt{instruction}} --- the natural language user query.
  \item \textbf{\texttt{tools}} --- JSON list of available tool definitions visible to the model.
  \item \textbf{\texttt{tool\_calls}} --- ground-truth list of \texttt{\{name, arguments\}} dicts.
  \item \textbf{\texttt{category}} --- tool domain label (e.g., \texttt{weather}, \texttt{stock}).
  \item \textbf{\texttt{num\_steps}} --- number of tool calls required (1 for single-step).
  \item \textbf{\texttt{text}} --- full formatted example in the target output format.
\end{itemize}

The canonical output format for training uses \texttt{<tool\_call>} XML tags:

\begin{lstlisting}
[weather] USER: What is the weather like in Tokyo?
ASSISTANT: <tool_call>
{"name": "get_weather", "arguments": {"city": "Tokyo", "units": "C"}}
</tool_call>
\end{lstlisting}

\subsection{Preprocessing Pipeline}

Raw JSONL examples are processed by \texttt{src/data\_loader.py} through
\texttt{ToolUseDataLoader}. The preprocessing steps are:

\begin{enumerate}[leftmargin=1.5em, itemsep=5pt]
  \item \textbf{Deduplication} --- MD5 hash of the full serialised row;
    first occurrence is kept.
  \item \textbf{Incomplete filtering} --- examples lacking any meaningful
    content field are dropped.
  \item \textbf{Whitespace normalisation} --- leading/trailing spaces removed.
  \item \textbf{Source balancing} --- enabled by default when multiple sources
    are configured.
  \item \textbf{Train/validation/test splitting} --- 80/10/10 ratio, seeded
    with 42 for reproducibility.
  \item \textbf{Tokenization} --- produces \texttt{input\_ids},
    \texttt{attention\_mask}, and \texttt{labels} columns; saved as
    Hugging Face Arrow files under \texttt{data/processed/}.
\end{enumerate}

A separate raw test split is saved to \texttt{data/processed/test\_raw.jsonl}
\textit{before} tokenization so that \texttt{src/evaluate.py} can access the
original prompt text, expected tool calls, and tool context metadata.

\subsection{Dataset Statistics}

\begin{table}[H]
\centering
\caption{Dataset statistics after preprocessing (default configuration).}
\begin{tabularx}{\linewidth}{L{5cm} C{9.2cm}}
\toprule
\thead{Property} & \thead{Value} \\
\midrule
Raw generated examples & 15\,000 (cap from \texttt{dataset\_config.yaml}) \\
After MD5 deduplication & ${\approx}$12\,280 \\
Training split ($80\%$) & ${\approx}$9\,824 \\
Validation split ($10\%$) & ${\approx}$1\,228 \\
Test split ($10\%$) & ${\approx}$1\,228 \\
Max sequence length & 2\,048 tokens \\
Language & English only \\
\bottomrule
\end{tabularx}
\end{table}

The loader reads all JSONL rows before applying the \texttt{samples=15000}
cap; the extra rows represent raw generator output that is subsequently
filtered, so the final split counts are as stated above.

\newpage

% =============================================================================
\section{System Prompt and Output Schema}
% =============================================================================

A shared system prompt defined in \texttt{src/constants.py} as
\texttt{TOOL\_USE\_SYSTEM\_PROMPT} is prepended to every conversation during
both training and evaluation. It specifies the exact output format the model
must produce:

\begin{lstlisting}
You are a helpful assistant that can use tools to answer user queries.

IMPORTANT INSTRUCTIONS:
1. Analyze the user's request and determine ALL tools needed to fully complete it.
2. If the request requires multiple steps, emit one <tool_call> block per step
   in the correct order.
3. Respond ONLY with tool call(s) - no explanations, no thinking text, no other output.
4. Each tool call must use <tool_call> and </tool_call> tags wrapping a JSON object
   with exactly 'name' and 'arguments' keys.

SINGLE-STEP FORMAT:
<tool_call>
{"name": "tool_name", "arguments": {"param_key": "param_value"}}
</tool_call>

MULTI-STEP FORMAT (emit all calls in sequence):
<tool_call>
{"name": "first_tool", "arguments": {"param": "value"}}
</tool_call>
<tool_call>
{"name": "second_tool", "arguments": {"param": "value"}}
</tool_call>
\end{lstlisting}

The schema requires \textbf{exactly two top-level keys} (\texttt{name} and
\texttt{arguments}) and prohibits any surrounding prose or chain-of-thought
text in the output. This strict format is what the reward functions measure
against.

\newpage

% =============================================================================
\section{Training Pipeline}
% =============================================================================

\subsection{Stage 1 --- Supervised Fine-Tuning}

\subsubsection{Objective}

SFT trains LoRA adapters on top of the frozen \texttt{Qwen/Qwen3-8B} weights
using a standard cross-entropy (next-token prediction) loss over synthetic
tool-use conversations. The goal is to shift the model's prior toward the
\texttt{<tool\_call>} output format before GRPO further refines compliance.

\subsubsection{Data Preparation for SFT}

Synthetic examples are loaded from \texttt{data/raw/synthetic/} by
\texttt{load\_synthetic\_conversations()} in \texttt{src/train.py}. Each
example is converted to a three-turn chat conversation:

\begin{enumerate}[leftmargin=1.5em, itemsep=3pt]
  \item \textbf{System} turn --- the shared \texttt{TOOL\_USE\_SYSTEM\_PROMPT}
    extended with the JSON schema of available tools for this example.
  \item \textbf{User} turn --- the natural language instruction.
  \item \textbf{Assistant} turn --- the ground-truth tool call(s) wrapped in
    \texttt{<tool\_call>} tags.
\end{enumerate}

Conversations are rendered into Tinker datums via the model's recommended
renderer (selected by \texttt{get\_recommended\_renderer\_name} from
\texttt{tinker\_cookbook}).

\subsubsection{LoRA Configuration}

Low-Rank Adaptation (LoRA) is used for parameter-efficient fine-tuning. Only
the adapter matrices $A$ and $B$ (where $\Delta W = BA$) are trained; the base
model weights remain frozen. The full LoRA configuration used in this pipeline
is:

\begin{table}[H]
\centering
\caption{LoRA adapter configuration for SFT.}
\begin{tabularx}{\linewidth}{L{5cm} C{9.2cm}}
\toprule
\thead{Hyperparameter} & \thead{Value} \\
\midrule
LoRA rank ($r$) & 64 \\
LoRA alpha ($\alpha$) & 16 \\
Target modules & \texttt{q\_proj}, \texttt{v\_proj}, \texttt{k\_proj}, \texttt{o\_proj}, \texttt{gate\_proj}, \texttt{up\_proj}, \texttt{down\_proj} \\
Adapter initialisation & $B = 0$, $A \sim \mathcal{N}(0, \sigma^2)$ (standard LoRA) \\
\bottomrule
\end{tabularx}
\end{table}

\subsubsection{SFT Hyperparameters}

\begin{table}[H]
\centering
\caption{SFT training hyperparameters (canonical pipeline defaults).}
\begin{tabularx}{\linewidth}{L{5cm} C{9.2cm}}
\toprule
\thead{Hyperparameter} & \thead{Value} \\
\midrule
Base model & \texttt{Qwen/Qwen3-8B} \\
LoRA rank & 64 \\
Learning rate & $2 \times 10^{-4}$ \\
Batch size & 8 \\
Epochs & 3 \\
Max sequence length & 2\,048 tokens \\
Optimizer & \texttt{paged\_adamw\_8bit} \\
Logging steps & 10 \\
Checkpoint save steps & 500 \\
Seed & 42 \\
Training platform & Tinker remote GPUs \\
\bottomrule
\end{tabularx}
\end{table}

\subsubsection{Checkpointing and Resumption}

Each SFT epoch writes a checkpoint entry to \texttt{outputs/sft/checkpoints.jsonl}
containing the Tinker sampler path. If training is interrupted, the script can
resume from the last valid checkpoint by reloading both model weights and
optimizer state (\texttt{load\_state\_with\_optimizer\_async}). The GRPO stage
reads the final sampler path from this file to initialise its starting checkpoint.

\subsubsection{Dry-Run Behaviour}

With \texttt{LOCAL\_VALIDATE=true}, the SFT script replaces remote Tinker calls
with a mock loop. The mock increments a simulated loss from 1.0 toward 0.0 at
$-0.05$ per step and logs W\&B metrics, allowing end-to-end pipeline wiring
to be verified without remote GPU access.

\vspace{8pt}

\subsection{Stage 2 --- Group Relative Policy Optimization}

\subsubsection{Algorithm}

GRPO is used as the reinforcement learning stage following SFT. It requires no
separate reward model and no human preference labels. For each training prompt,
the model samples a \textbf{group of $G$ completions}. Each completion is scored
by the programmatic reward functions in \texttt{src/rewards.py}. The
advantage of completion $i$ within its group is:

\begin{equationbox}
\[
  \hat{A}_i = \frac{r_i - \mathrm{mean}(r_{1\ldots G})}{\mathrm{std}(r_{1\ldots G})}
\]
\end{equationbox}

This normalised advantage is used to update the policy via a PPO-style clipped
surrogate objective. Groups where all completions receive the same reward
(degenerate groups) provide no learning signal and are skipped. The metric
\texttt{train/frac\_degenerate} tracks the proportion of such groups per step.

\subsubsection{GRPO Hyperparameters}

\begin{table}[H]
\centering
\caption{GRPO training hyperparameters (canonical pipeline defaults).}
\begin{tabularx}{\linewidth}{L{5cm} C{9.2cm}}
\toprule
\thead{Hyperparameter} & \thead{Value} \\
\midrule
Base model & \texttt{Qwen/Qwen3-8B} \\
Starting checkpoint & Latest valid checkpoint from \texttt{outputs/sft/} \\
LoRA rank & 64 (must match SFT) \\
Learning rate & $4 \times 10^{-5}$ \\
Batch size & 16 \\
Group size ($G$) & 8 \\
Max training steps & 50 \\
Max completion length & 512 tokens \\
Checkpoint save steps & 17 \\
Seed & 42 \\
Training platform & Tinker remote GPUs \\
\bottomrule
\end{tabularx}
\end{table}

The requirement that the GRPO LoRA rank matches the SFT rank is enforced in
\texttt{scripts/run\_pipeline.sh}; mismatched ranks cause an immediate
pipeline abort before any Tinker API calls.

\subsubsection{Prompt Dataset for GRPO}

GRPO prompts are loaded through \texttt{ToolUseDataLoader} in
\texttt{build\_prompt\_dataset()}. Each prompt dict carries:

\begin{itemize}[leftmargin=1.5em, itemsep=3pt]
  \item \textbf{\texttt{prompt\_text}} --- the natural language instruction.
  \item \textbf{\texttt{expected\_tool}} --- the correct first tool name (for
    \texttt{tool\_name\_reward}).
  \item \textbf{\texttt{expected\_args}} --- the correct argument dict (for
    \texttt{argument\_f1\_reward}).
  \item \textbf{\texttt{expected\_chain}} --- the ordered list of tool names
    for multi-step examples (for chain rewards).
  \item \textbf{\texttt{source}} --- the originating data source tag.
\end{itemize}

\subsubsection{Dry-Run Behaviour}

With \texttt{LOCAL\_VALIDATE=true}, the GRPO script logs simulated reward and
degenerate-group metrics through the same W\&B integration used in full training,
allowing the reward computation and group-update logic to be verified locally.

% =============================================================================
\section{Reward Function Design}
% =============================================================================

The reward functions in \texttt{src/rewards.py} are the core of the GRPO stage.
All functions share the same call signature:
\texttt{reward\_fn(completions, **kwargs) -> list[float]}.

\subsection{Tool Call Extraction}

Before scoring, completions are parsed by \texttt{extract\_tool\_call()} (single
call) or \texttt{extract\_tool\_calls()} (all calls). Both functions first strip
any \texttt{<think>...</think>} blocks (for reasoning-enabled models), then
attempt extraction in priority order:

\begin{enumerate}[leftmargin=1.5em, itemsep=3pt]
  \item Fenced JSON blocks: \texttt{```json \{...\} ```}
  \item \texttt{<tool\_call>...</tool\_call>} XML tags (primary format)
  \item Qwen \texttt{<|function\_call|>} format
  \item Fallback: any JSON object with a \texttt{"name"} key
\end{enumerate}

\subsection{Individual Reward Signals}

\begin{table}[H]
\centering
\caption{Reward functions and their behaviours.}
\begin{tabularx}{\linewidth}{L{4.5cm} C{1.5cm} L{8.2cm}}
\toprule
\thead{Reward Function} & \thead{Range} & \thead{Description} \\
\midrule
\texttt{schema\_validation\_reward} & $[0, 1]$ & Binary. 1.0 if the completion
  contains a parseable tool call with a non-empty \texttt{name} string and a
  dict \texttt{arguments} field; 0.0 otherwise. Always applied. \\[4pt]
\texttt{tool\_name\_reward} & $[0, 1]$ & Binary. 1.0 if the predicted tool name
  matches \texttt{expected\_tool} (case-insensitive); 0.0 otherwise. Applied
  when \texttt{expected\_tool} is provided. \\[4pt]
\texttt{argument\_f1\_reward} & $[0, 1]$ & Partial credit. Computes key-value
  F1 between predicted and expected argument dicts (string-normalised values).
  Applied when \texttt{expected\_args} is provided. \\[4pt]
\texttt{full\_chain\_reward} & $[0, 1]$ & Binary. 1.0 if the \textit{ordered}
  sequence of predicted tool names exactly matches \texttt{expected\_chain}.
  Applied for multi-step examples. \\[4pt]
\texttt{chain\_partial\_reward} & $[0, 1]$ & Partial credit. Position-wise
  fraction of expected chain tools correctly predicted ($\frac{1}{N}$ per
  correct position). Provides gradient signal for partially correct chains. \\
\bottomrule
\end{tabularx}
\end{table}

\subsection{Composite Reward}

\texttt{compute\_rewards()} in \texttt{src/rewards.py} assembles the applicable
reward functions at runtime based on which metadata fields are available in the
prompt dict. The composite reward is the \textbf{arithmetic mean} of all
applicable component rewards:

\begin{equationbox}
\[
  r_{\text{composite}} = \frac{1}{|\mathcal{F}|}\sum_{f \in \mathcal{F}} r_f
\]
\end{equationbox}

where $\mathcal{F}$ is the set of active reward functions for the current
example. This equal-weight averaging ensures that all applicable signals
contribute proportionally and avoids any single signal dominating.

The mixing strategy deliberately combines \textbf{strict binary signals}
(schema, tool name, full chain) with \textbf{partial-credit signals}
(argument F1, chain partial). Binary signals establish the categorical
constraint; partial-credit signals provide smooth gradients when the model
is near but not exactly correct.

% =============================================================================
\section{Evaluation Framework}
% =============================================================================

\subsection{Evaluation Harness}

\texttt{src/evaluate.py} implements a three-stage evaluation harness that
can compare baseline, SFT, and GRPO checkpoints in a single run. Evaluation
is driven by Tinker inference (remote model sampling) against the structured
test split.

\subsection{Test Data Source}

Evaluation examples are loaded in priority order:

\begin{enumerate}[leftmargin=1.5em, itemsep=3pt]
  \item \texttt{data/processed/test\_raw.jsonl} --- preferred; preserves tool
    metadata, prompt text, and structured expected fields.
  \item Held-out examples from \texttt{data/raw/synthetic/} --- used if the
    prepared test file is absent; a stable 10\% hash-based split ensures the
    same examples are always selected.
\end{enumerate}

A stable test split is implemented via MD5 hashing of each row's serialised
JSON; rows whose hash modulo 100 falls below the threshold (10) are included.
This ensures reproducibility without a fixed ordering requirement.

\subsection{Evaluation Metrics}

\begin{table}[H]
\centering
\caption{Evaluation metrics computed for each stage.}
\begin{tabularx}{\linewidth}{L{5cm} L{9.2cm}}
\toprule
\thead{Metric Key} & \thead{Definition} \\
\midrule
\texttt{tool\_selection\_accuracy} & Exact match on the first extracted tool call name. \\
\texttt{argument\_accuracy} & Exact dictionary match on extracted arguments. \\
\texttt{schema\_compliance} & Output contains a parseable tool call with valid name and argument dict. \\
\texttt{multi\_step\_success} & Expected tool-call chain reproduced in correct order. \\
\texttt{multi\_step\_partial} & At least one but not all chain steps correct. \\
\texttt{multi\_step\_one\_call\_only} & Only one call produced for a multi-step prompt. \\
\texttt{multi\_step\_wrong\_order} & Correct tools but in wrong order. \\
\texttt{multi\_step\_zero\_calls} & No tool call produced for a multi-step prompt. \\
\texttt{avg\_latency\_ms} & Average wall-clock time per model response (milliseconds). \\
\texttt{ms\_per\_token} & Normalised latency per generated token. \\
\bottomrule
\end{tabularx}
\end{table}

\subsection{Pipeline Commands}

\begin{table}[H]
\centering
\caption{Evaluation entrypoints from \texttt{src/evaluate.py} and the pipeline script.}
\begin{tabularx}{\linewidth}{L{4cm} L{10.2cm}}
\toprule
\thead{Mode} & \thead{Command / Purpose} \\
\midrule
\texttt{baseline} & Zero-shot evaluation of \texttt{Qwen/Qwen3-8B} with no adapter \\
\texttt{sft} & Evaluate the SFT checkpoint \\
\texttt{grpo} & Evaluate the GRPO checkpoint \\
\texttt{compare} & Run all three stages and produce a merged JSON \\
\texttt{all} & Full pipeline (baseline + sft + grpo + compare) \\
\texttt{local-compare} & Reconstruct comparison from pre-saved JSON result files \\
\texttt{fetch-compare} & Reconstruct comparison from W\&B run data \\
\bottomrule
\end{tabularx}
\end{table}

\newpage

% =============================================================================
\section{Results and Analysis}
% =============================================================================

\subsection{Cross-Stage Comparison}

The canonical results are stored in \texttt{outputs/eval\_comparison.json}
and reproduced in full below.

\begin{table}[H]
\centering
\caption{Full evaluation results across all three training stages.}
\begin{tabularx}{\linewidth}{L{5.5cm} C{2.5cm} C{2.5cm} C{3.7cm}}
\toprule
\thead{Metric} & \thead{Baseline} & \thead{SFT} & \thead{GRPO} \\
\midrule
Tool Selection Accuracy       & 0.782 & 0.774 & 0.774 \\
Argument Accuracy             & 0.308 & 0.310 & \textbf{0.321} \\
Schema Compliance             & 0.807 & 0.812 & \textbf{0.814} \\
Multi-Step Success            & 0.406 & 0.386 & \textbf{0.413} \\
Multi-Step Partial            & 0.394 & 0.402 & \textbf{0.425} \\
Multi-Step One-Call Only      & 0.039 & 0.031 & 0.035 \\
Multi-Step Wrong Order        & 0.004 & 0.012 & 0.020 \\
Multi-Step Zero Calls         & 0.154 & \textbf{0.165} & \textbf{0.106} \\
Avg Latency (ms)              & 6\,074.8 & 7\,124.9 & 7\,098.9 \\
ms per Token                  & 15.79 & 20.73 & 20.42 \\
\bottomrule
\end{tabularx}
\end{table}

\subsection{Key Observations}

\subsubsection{Tool Selection Accuracy}

The baseline \texttt{Qwen3-8B} achieves 78.2\% tool selection accuracy
zero-shot, reflecting the strong pre-training of the Qwen3 series. Both SFT and
GRPO register 77.4\%, a marginal decline. This is consistent with the model
having been trained on synthetic data whose instruction distribution is narrower
than the zero-shot capability --- fine-tuning on a constrained domain can
slightly regress generalisation while improving schema-specific behaviour.

\subsubsection{Argument Accuracy}

Argument accuracy shows a monotonic improvement across stages: 0.308
(baseline) $\rightarrow$ 0.310 (SFT) $\rightarrow$ 0.321 (GRPO). The
\texttt{argument\_f1\_reward} applied during GRPO provides a direct partial-credit
gradient signal for argument key--value correctness, explaining this improvement.
The 1.3 percentage point gain from SFT to GRPO is modest, reflecting the
limited number of GRPO steps (50) relative to the diversity of argument patterns
in the dataset.

\subsubsection{Schema Compliance}

Schema compliance rises from 80.7\% (baseline) to 81.2\% (SFT) to 81.4\%
(GRPO). The baseline's 80.7\% is already high because \texttt{Qwen3-8B} natively
understands structured output. The marginal gains confirm that the \texttt{schema\_validation\_reward}
in GRPO successfully reinforces the \texttt{<tool\_call>} output format.

\subsubsection{Multi-Step Performance}

Multi-step metrics reveal the most interesting pattern:

\begin{itemize}[leftmargin=1.5em, itemsep=5pt]
  \item \textbf{Multi-Step Success} drops from 0.406 (baseline) to 0.386
    (SFT), then recovers to 0.413 (GRPO). This U-shape is important: SFT on
    single-step-dominated synthetic data may cause the model to regress on
    multi-step problems. GRPO reverses this regression and surpasses the
    baseline.
  \item \textbf{Multi-Step Zero Calls} --- the fraction of multi-step prompts
    where the model emits \textit{no} tool call at all --- falls substantially
    from 16.5\% (SFT) to 10.6\% (GRPO). This is the clearest signature of
    GRPO's reinforcement signal working: the \texttt{schema\_validation\_reward}
    penalises zero-call outputs directly.
  \item \textbf{Multi-Step Partial} improves from 0.402 (SFT) to 0.425 (GRPO),
    indicating that the model is producing more of the required chain steps even
    when it does not complete all of them.
  \item \textbf{Multi-Step Wrong Order} increases slightly (0.012 $\rightarrow$
    0.020), suggesting that GRPO pushes the model to attempt more multi-step
    outputs, and some of those attempts get the order wrong. This is an expected
    trade-off with more aggressive exploration.
\end{itemize}

\subsubsection{Latency}

Both fine-tuned stages are slower than the zero-shot baseline
(7\,100--7\,125\,ms vs. 6\,075\,ms average latency). This increase is expected:
the LoRA adapters add computation to each forward pass, and the longer outputs
required by the \texttt{<tool\_call>} format increase token count relative to
bare natural-language responses. GRPO (7\,099\,ms) is marginally faster than
SFT (7\,125\,ms) on average, with a lower ms-per-token (20.42 vs. 20.73).

\subsection{Intermediate Test Files}

The repository also contains three lightweight intermediate test files in
\texttt{outputs/}:

\begin{table}[H]
\centering
\caption{Intermediate single-metric evaluation snapshots.}
\begin{tabularx}{\linewidth}{L{5.5cm} C{8.7cm}}
\toprule
\thead{File} & \thead{Tool Selection Accuracy} \\
\midrule
\texttt{\_tmp\_eval\_baseline.json} & 0.50 \\
\texttt{\_tmp\_eval\_sft.json} & 0.70 \\
\texttt{\_tmp\_eval\_grpo.json} & 0.90 \\
\texttt{\_tmp\_eval\_compare.json} & baseline 0.50, SFT 0.70, GRPO 0.90 \\
\bottomrule
\end{tabularx}
\end{table}

These files contain only a single metric (\texttt{tool\_selection\_accuracy})
and appear to be lightweight test runs or placeholder outputs used to verify
pipeline wiring. They are not
representative of the full evaluation; \texttt{outputs/eval\_comparison.json}
is the authoritative multi-metric result.



\newpage

% =============================================================================
\section{Pipeline Execution and Reproducibility}
% =============================================================================

\subsection{End-to-End Entrypoint}

The canonical pipeline is invoked via a single shell script:

\begin{lstlisting}
bash scripts/run_pipeline.sh all
\end{lstlisting}

Stage-specific invocations are also available:

\begin{lstlisting}
bash scripts/run_pipeline.sh baseline   # zero-shot eval only
bash scripts/run_pipeline.sh sft        # SFT training + eval
bash scripts/run_pipeline.sh grpo       # GRPO training + eval
bash scripts/run_pipeline.sh compare    # cross-stage comparison
\end{lstlisting}

\subsection{Local Validation Mode}

To validate pipeline wiring without remote compute:

\begin{lstlisting}
LOCAL_VALIDATE=true bash scripts/run_pipeline.sh all
\end{lstlisting}

This activates the following substitutions:
\begin{itemize}[leftmargin=1.5em, itemsep=3pt]
  \item Base model: \texttt{Qwen/Qwen3-8B} $\rightarrow$ \texttt{sshleifer/tiny-gpt2}
  \item Training: remote Tinker calls $\rightarrow$ local mock loops
  \item Hub upload: disabled (\texttt{HF\_REPO\_ID} cleared)
  \item Dry-run steps: 3 (configurable via \texttt{--dry-run-steps})
\end{itemize}

\subsection{Pre-Flight Checks}

The pipeline script performs pre-flight validation before any remote API call:

\begin{itemize}[leftmargin=1.5em, itemsep=3pt]
  \item \texttt{TINKER\_API\_KEY} must be set (unless \texttt{LOCAL\_VALIDATE=true}).
  \item The synthetic data directory must contain at least one JSONL file
    (\texttt{data/raw/synthetic}); if not, the generator script
    (\texttt{scripts/generate\_synthetic.py}) is invoked automatically.
  \item GRPO LoRA rank must match the SFT LoRA rank; mismatch causes immediate
    abort with a descriptive error.
  \item Warnings are emitted (but execution continues) when \texttt{WANDB\_API\_KEY}
    is absent, or when \texttt{HF\_REPO\_ID} is set without \texttt{HF\_TOKEN}.
\end{itemize}

\subsection{Optional Hugging Face Hub Upload}

Both the SFT and GRPO scripts support uploading the final LoRA adapter to the
Hugging Face Hub via \texttt{push\_model\_to\_hub.py}. This is activated by
setting \texttt{HF\_TOKEN} and \texttt{HF\_REPO\_ID} in the environment.

% =============================================================================
\section{Limitations and Improvement Opportunities}
% =============================================================================

This repository demonstrates a robust end-to-end training and evaluation
pipeline, but several limitations are visible from the current implementation
and outputs.

\subsection{Current Limitations}

\begin{enumerate}[leftmargin=1.5em, itemsep=6pt]
  \item \textbf{Synthetic-only training distribution.}
    All training data is generated from templates. This ensures controllability
    and reproducibility, but limits linguistic diversity and real-world tool-use
    variability. The model may overfit to templated phrasing and underperform on
    naturally noisy user requests.

  \item \textbf{Narrow evaluation domain.}
    Evaluation examples are sourced from the prepared processed test split
    (\texttt{test\_raw.jsonl} under \texttt{data/processed}) or held-out
    synthetic rows, so test data remains close to the training generation
    process. This likely inflates in-domain metrics relative to real-world
    prompts.

  \item \textbf{Modest GRPO training budget.}
    The canonical GRPO run uses 50 steps with group size 8. This is sufficient
    to show measurable gains, but small for a model of this scale and may not
    fully converge reward-conditioned behavior.

  \item \textbf{Multi-step ordering errors persist.}
    \texttt{multi\_step\_wrong\_order} increases from 0.012 (SFT) to 0.020 (GRPO),
    indicating that while GRPO encourages more chain attempts, ordering
    consistency remains an unresolved failure mode.

  \item \textbf{Latency overhead after fine-tuning.}
    SFT and GRPO checkpoints are slower than baseline (around 7.1s vs 6.1s avg
    latency). For interactive agents, this can materially affect user
    experience.

  \item \textbf{No production-serving benchmark in repository.}
    The current results quantify inference latency in evaluation mode but do not
    include throughput/load tests under concurrent serving settings.
\end{enumerate}

\subsection{Improvement Opportunities}

\begin{enumerate}[leftmargin=1.5em, itemsep=6pt]
  \item \textbf{Blend synthetic and real traces.}
    Add curated real conversation logs (with privacy filtering) or open
    function-calling corpora to reduce distribution shift and improve
    generalization.

  \item \textbf{Strengthen multi-step reward shaping.}
    Add explicit penalties for wrong-order chains and repeated calls, and add a
    positive reward for terminating with a concise final answer after successful
    tool execution.

  \item \textbf{Increase GRPO training depth.}
    Run longer GRPO schedules (for example 150--400 steps), monitor per-component
    reward curves, and evaluate whether gains saturate or continue.

  \item \textbf{Expand evaluation benchmarks.}
    Introduce adversarial and out-of-domain prompt suites (ambiguous queries,
    missing arguments, conflicting constraints, and malformed tool specs).

  \item \textbf{Add schema-aware decoding constraints.}
    Use constrained decoding or grammar-guided decoding during inference to
    reduce malformed outputs and improve first-pass compliance.

  \item \textbf{Improve serving efficiency.}
    Introduce quantized inference checkpoints and KV-cache-aware serving
    benchmarks to recover latency while preserving structured-output quality.
\end{enumerate}

\newpage

% =============================================================================
\section{Conclusions}
% =============================================================================

Based on the implementation and results of this repository, the following
conclusions can be drawn

\begin{enumerate}[leftmargin=1.5em, itemsep=10pt]

  \item \textbf{SFT on synthetic data provides incremental gains on argument
    accuracy and schema compliance, but can regress multi-step performance.}
    SFT raised argument accuracy from 0.308 to 0.310 and schema compliance
    from 0.807 to 0.812. However, multi-step success dropped from 0.406 to
    0.386, suggesting that single-step-dominated synthetic data partially
    suppresses the base model's multi-step reasoning capability.

  \item \textbf{GRPO recovers and surpasses the baseline on multi-step metrics.}
    GRPO pushed multi-step success to 0.413 (above the 0.406 baseline) and
    reduced the zero-call failure rate from 16.5\% to 10.6\%. The
    \texttt{schema\_validation\_reward} directly penalises empty outputs,
    providing an explicit gradient signal that SFT's cross-entropy loss cannot
    supply.

  \item \textbf{Argument quality improves most under GRPO.}
    The \texttt{argument\_f1\_reward} provides a partial-credit signal for
    argument key--value overlap, which SFT's binary cross-entropy cannot
    express. This explains why argument accuracy improves more from SFT to GRPO
    (0.310 $\rightarrow$ 0.321) than from baseline to SFT (0.308 $\rightarrow$
    0.310).

  \item \textbf{The composite, multi-signal reward design is effective.}
    By combining strict binary signals (schema validation, exact tool name,
    exact chain match) with partial-credit signals (argument F1, positional
    chain credit), the reward function provides both categorical constraints
    and smooth gradients.

  \item \textbf{Fine-tuning introduces a latency cost.}
    Both SFT and GRPO checkpoints are approximately 17\% slower than the
    zero-shot baseline (7\,100\,ms vs. 6\,075\,ms average). This overhead
    is attributable to LoRA adapter computation and longer structured outputs.
    GRPO is marginally faster than SFT in per-token terms (20.42 vs.
    20.73\,ms/token).

  \item \textbf{The pipeline is fully reproducible and infrastructure-agnostic.}
    The dry-run mode (\texttt{LOCAL\_VALIDATE=true}) allows complete
    end-to-end pipeline validation without remote GPU access. The
    synthetic data generator uses a fixed seed, preprocessing is deterministic,
    and all hyperparameters are explicitly versioned in the pipeline script.

\end{enumerate}

\newpage

% =============================================================================
\section*{References}
% =============================================================================
\addcontentsline{toc}{section}{References}

\begin{enumerate}[leftmargin=1.5em, itemsep=6pt, label={[\arabic*]}]

  \item Qwen Team, Alibaba Cloud. (2025). Qwen3 Technical Report.
    \url{https://huggingface.co/Qwen/Qwen3-8B}

  \item Hu, E.\,J., Shen, Y., Wallis, P., Allen-Zhu, Z., Li, Y., Wang, S.,
    Wang, L., \& Chen, W. (2021). LoRA: Low-Rank Adaptation of Large Language
    Models. \textit{arXiv preprint arXiv:2106.09685}.

  \item DeepSeek-AI. (2025). DeepSeek-R1: Incentivizing Reasoning Capability in
    LLMs via Reinforcement Learning. \textit{arXiv preprint arXiv:2501.12948}.

  \item Schulman, J., Wolski, F., Dhariwal, P., Radford, A., \& Klimov, O. (2017).
    Proximal Policy Optimization Algorithms. \textit{arXiv preprint
    arXiv:1707.06347}.

  \item Rafailov, R., Sharma, A., Mitchell, E., Manning, C.\,D., Ermon, S., \&
    Finn, C. (2023). Direct Preference Optimization: Your Language Model is
    Secretly a Reward Model. \textit{arXiv preprint arXiv:2305.18290}.

  \item Ouyang, L., et al. (2022). Training language models to follow instructions
    with human feedback. \textit{arXiv preprint arXiv:2203.02155}.

  \item Tinker --- Remote GPU Infrastructure for LLM Training and Inference.
    \url{https://tinker.thinkingmachines.ai/}

  \item Wandb Inc. (2020). Weights \& Biases: Machine Learning Experiment
    Tracking. \url{https://wandb.ai}

\end{enumerate}

\end{document}
